\section{Discussion and open questions}

A score label linked to each trial row preserves the association between a measurable set of assignment probabilities, treatment, and outcome. Given a known score law, the arm-label outcome and score distributions determine the sharp mean endpoints in \cref{obj:def:mean-endpoints} and hence the average treatment effect interval in \cref{obj:def:sharp-ate-set}. The worst-case width in \cref{obj:def:worst-case-ambiguity} measures the information retained by a release in units of the target parameter. \Cref{obj:thm:universal-point-identification} characterizes when that information suffices for point identification across compatible released laws.

For score laws on \([\varepsilon,1-\varepsilon]\) with a common density floor \(m_f\) satisfying \(0<m_f\le(1-2\varepsilon)^{-1}\), \cref{obj:thm:k-label-rate} gives an order-\(K^{-1}\) rate for optimal worst-case ambiguity. Under the score-law conditions of \cref{obj:thm:sharp-high-resolution-constant}, its high-resolution allocation rule assigns finer cells where local variation contributes more to the ambiguity criterion. The rule connects the placement of labels to the sharp interval that an analyst can recover from the released data. \Cref{obj:thm:minimax-honest-length} establishes a \(K^{-1}+n^{-1/2}\) minimax honest-length rate for known score laws with a positive density lower bound, attained up to constants by equal-width labels.

\subsection{Open questions}

The projection interval in \cref{obj:def:external-projection-handle} is defined through a countable rational sieve. A finite truncation and optimization rule, with an error bound that preserves coverage and controls excess length for atomic laws and arbitrary measurable cells, remains to be developed.

The two-source bounds in \cref{obj:thm:external-score-root-order} establish the root scale for excess interval length with an independent score log. A sharp limit would quantify the combined contribution of trial and score-log sampling when their sizes grow at a fixed positive ratio. The question is stated for the full external-law class and, where appropriate, for a specified regular subclass.

\begin{remarkv}{oeq:external-score-projection}[External-score sharp limits]
Fix a measurable release \(g\), \(0<\alpha<1/2\), and \(\rho\in(0,\infty)\). The sharp asymptotic analysis over \(\mathfrak X(g)\), with confidence procedures in \(\mathcal J^{\mathrm{ext}}_{n,m,\alpha}(g)\), concerns
\[
\liminf_{\substack{n,m\to\infty\\ n/m\to\rho}}T_{n,m,\alpha}(g)
\qquad\text{and}\qquad
\limsup_{\substack{n,m\to\infty\\ n/m\to\rho}}T_{n,m,\alpha}(g).
\]
It asks whether these limits coincide and, if so, for their common sharp constant. If a common sharp limit fails over the full class, the corresponding question concerns an explicitly defined regular subclass of \(\mathfrak X(g)\), with both uniform honesty and the supremum in expected excess length restricted to that subclass. The subclass analysis specifies its regularity conditions and sharp limit for each \(\rho\in(0,\infty)\), using the external projection construction to obtain an attaining procedure or a counterexample.
\end{remarkv}

For a regular subclass, the honesty requirement and worst-case excess-length criterion in \cref{obj:oeq:external-score-projection} apply to the same collection of laws. The projection construction in \cref{obj:def:external-projection-handle} supplies a concrete route for studying attainment, while the fixed sample-size ratio keeps both sources of sampling variation in the asymptotic comparison.
